\section*{अध्याय \ref{ch:g10:proba} --- प्रायिकता और प्रतिदर्शन}

\begin{solution}{exo:g10:proba:1}
$6$ परिणामों में से हर एक की \omterm{def:g10:proba:distribution}{प्रायिकता} $\frac16$ है।
``$6$ आना'': $\frac16$।
``विषम संख्या'' $= \{1, 3, 5\}$: $\frac36 = \frac12$।
``अधिक से अधिक $4$'' $= \{1,2,3,4\}$: $\frac46 = \frac23$।
``$7$ आना'': असंभव \omterm{def:g10:proba:events}{घटना}, \omterm{def:g10:proba:distribution}{प्रायिकता} $0$।
\end{solution}

\begin{solution}{exo:g10:proba:2}
कुल $10$ समसंभावी कंचे हैं।
$\P(\text{हरा}) = \frac{5}{10} = \frac12$।
$\P(\text{काला नहीं}) = 1 - \frac{2}{10} = \frac{8}{10} = \frac45$।
हरा और पीला \omterm{def:g10:proba:operations}{असंगत} हैं: $\P(\text{हरा या पीला}) = \frac{5}{10} + \frac{3}{10} = \frac{8}{10}
= \frac45$।
\end{solution}

\begin{solution}{exo:g10:proba:3}
मान लीजिए $S$: ``स्पेनी पढ़ता है'', $G$: ``जर्मन पढ़ता है''। तब $\P(S) = \frac{18}{30}$,
$\P(G) = \frac{10}{30}$, $\P(S \cap G) = \frac{4}{30}$। योग का नियम:
\[
\P(S \cup G) = \frac{18}{30} + \frac{10}{30} - \frac{4}{30}
= \frac{24}{30} = \frac45 .
\]
\omterm{def:g10:proba:operations}{पूरक}: $\P(\text{कोई नहीं}) = 1 - \frac45 = \frac15$।
\end{solution}

\begin{solution}{exo:g10:proba:4}
\omterm{met:g10:proba:tree}{वृक्ष} के चार पथ HH, HT, TH, TT हैं, और हर एक की \omterm{def:g10:proba:distribution}{प्रायिकता} $\frac12 \times \frac12 = \frac14$ है।
दो चित: $\frac14$। ठीक एक चित (HT या TH): $\frac24 = \frac12$। कम से कम एक चित: $1 - \P(\text{TT}) = 1 - \frac14 = \frac34$।
\end{solution}

\begin{solution}{exo:g10:proba:5}
पहली बार: लाल \omterm{def:g10:proba:distribution}{प्रायिकता} $\frac{4}{10}$ से, सफ़ेद $\frac{6}{10}$ से।
दूसरी बार \emph{बिना} वापसी: लाल के बाद $9$ में से $3$ लाल और $6$ सफ़ेद
बचती हैं; सफ़ेद के बाद $9$ में से $4$ लाल और $5$ सफ़ेद।

दो लाल: $\frac{4}{10} \times \frac39 = \frac{12}{90} = \frac{2}{15}$।

अलग-अलग रंग: पथ RW और WR:
\[
\frac{4}{10} \times \frac69 + \frac{6}{10} \times \frac49
= \frac{24}{90} + \frac{24}{90} = \frac{48}{90} = \frac{8}{15}.
\]
\end{solution}

\begin{solution}{exo:g10:proba:6}
\emph{1.} हर पासा स्वतंत्र रूप से हर फलक \omterm{def:g10:proba:distribution}{प्रायिकता} $\frac16$ से दिखाता है, इसलिए
$36$ क्रमित युग्म समसंभावी हैं; और योग युग्म से तय हो जाता है।

\emph{2.} योग $7$: युग्म $(1,6), (2,5), (3,4), (4,3), (5,2), (6,1)$, इसलिए $\P = \frac{6}{36} = \frac16$। योग $12$: केवल $(6,6)$, $\P = \frac{1}{36}$। योग
अधिक से अधिक $4$: वे युग्म जिनका योग $2$, $3$ या $4$ है: $(1,1)$; $(1,2), (2,1)$;
$(1,3), (2,2), (3,1)$ --- छह युग्म, इसलिए $\P = \frac{6}{36} = \frac16$।

\emph{3.} हर योग को देने वाले युग्मों की संख्या $1$ (योग $2$) से बढ़कर
$6$ (योग $7$) तक जाती है और फिर घटती है: सबसे संभावित योग $7$ है।
\end{solution}

\begin{solution}{exo:g10:proba:7}
हर प्रश्न का उत्तर स्वतंत्र रूप से \omterm{def:g10:proba:distribution}{प्रायिकता} $\frac14$ से सही होता है। दो स्तरों
वाला \omterm{met:g10:proba:tree}{वृक्ष}:

दोनों सही: $\frac14 \times \frac14 = \frac{1}{16}$।

ठीक एक सही: पथ सही--ग़लत और ग़लत--सही: $\frac14 \times \frac34 + \frac34 \times \frac14 = \frac{6}{16}
= \frac38$।

एक भी सही नहीं: $\frac34 \times \frac34 = \frac{9}{16}$।

जाँच: $\frac{1}{16} + \frac{6}{16} + \frac{9}{16} = 1$।
\end{solution}

\begin{solution}{exo:g10:proba:8}
\emph{1.} प्रेक्षित \omterm{def:g10:stats:series}{आपेक्षिक बारंबारता}: $\frac{396}{900} = 0.44$। $n = 900$ के साथ: $\frac{1}{\sqrt{900}} = \frac{1}{30} \approx 0.033$, इसलिए
उतार-चढ़ाव \omterm{def:g10:numbers:interval}{अंतराल} लगभग $\intcc{0.467}{0.533}$ है।

\emph{2.} प्रेक्षित $0.44$ उस \omterm{def:g10:numbers:interval}{अंतराल} से बाहर है: $50\%$ समर्थन वाली समष्टि से
लिए गए $900$ मतदाताओं के \omterm{def:g10:proba:sample}{प्रतिदर्श} इतनी दूर विरले ही भटकते हैं। सर्वेक्षण
दावे किए गए $50\%$ पर गंभीर संदेह खड़ा करता है।
\end{solution}

\begin{solution}{exo:g10:proba:9}
\emph{1.} $\P(\text{जीत } 6) = \frac16$ ($6$ आने पर), $\P(\text{जीत } 3) = \frac16$ ($5$ आने पर), $\P(\text{जीत } 0) = \frac46 = \frac23$।

\emph{2.} $600$ खेलों में लगभग $100$ छक्के, $100$ पंजे और $400$ अन्य परिणाम
अपेक्षित हैं। जीत: लगभग $100 \times 6 + 100 \times 3 = 900$। लागत: $600 \times 2 = 1200$। अपेक्षित शुद्ध परिणाम: $900 - 1200 = -300$,
अर्थात् औसतन प्रति खेल लगभग $0.5$ की हानि --- यह खेल खेलने लायक़ नहीं है।
\end{solution}

\begin{solution}{pb:g10:proba:1}
\textbf{1.} $P(\text{बादशाह या पान}) = \frac{4}{52} +
\frac{13}{52} - \frac{1}{52} = \frac{16}{52} = \frac{4}{13}$: पान का बादशाह दोनों \omterm{def:g10:proba:events}{घटनाओं} में बैठा है और उसे दो बार नहीं
गिनना चाहिए --- इसीलिए घटाव (\cref{thm:g10:proba:addition})।

\textbf{2.} $P(\text{कोई चित नहीं}) = \left(\frac12\right)^3 =
\frac18$, इसलिए $P(\text{कम से कम एक}) = \frac78$।

\textbf{3.} $P(RR) = \frac35 \times \frac24 = \frac{3}{10}$; $P(\text{प्रत्येक का एक}) = \frac35 \times \frac24 + \frac25
\times \frac34 = \frac{6}{20} + \frac{6}{20} = \frac35$।

\textbf{4.} $P(BB) = \frac25 \times \frac14 = \frac1{10}$; और $\frac{3}{10} + \frac{6}{10} + \frac{1}{10} = 1$: \omterm{met:g10:proba:tree}{वृक्ष} के पत्तों का योग सदा $1$ होता है।

\textbf{5.} \omterm{def:g10:proba:operations}{असंगत} \omterm{def:g10:proba:events}{घटनाएँ} $P(A \cup B) = P(A) +
P(B) = 1.1 > 1$ मानती हैं: असंभव --- कहीं से \omterm{def:g10:proba:distribution}{प्रायिकता} गढ़
ली गई है।

\textbf{6.} (अधिकांश लोग, उन हज़ार पत्र-लेखकों सहित, कहते हैं ``कोई अंतर
नहीं पड़ता, 50--50''। अपना उत्तर प्रश्न 8 के लिए सँभालकर रखिए।)

\textbf{7.} कार दरवाज़े 1 के पीछे (\omterm{def:g10:proba:distribution}{प्रायिकता} $\frac13$): संचालक दोनों में से कोई
भी बकरी-दरवाज़ा खोल देता है; बदलने पर कार हाथ से निकल जाती है: \emph{हार}।
कार दरवाज़े 2 के पीछे ($\frac13$): संचालक को दरवाज़ा 3 खोलना \emph{ही} पड़ता है;
बदलने पर आप दरवाज़े 2 पर आ जाते हैं: \emph{जीत}। कार दरवाज़े 3 के पीछे:
सममित: \emph{जीत}। कुल: $P(\text{बदलने पर जीत}) = \frac13 + \frac13 = \frac23$।

\textbf{8.} आपका पहला चुनाव \omterm{def:g10:proba:distribution}{प्रायिकता} $\frac23$ से ग़लत होता है --- और ठीक तभी
संचालक की मजबूरी में खोली गई बकरी बची हुई दरवाज़े के पीछे कार छोड़ देती है:
बदलना हर आरंभिक भूल को जीत में बदल देता है। $\frac23$, चाहे पहला अनुमान कुछ भी
कहता रहा हो।

\textbf{9.} छह समसंभावी परिदृश्य: कार $1$, $2$ या $3$ के पीछे, और
संचालक दरवाज़ा $2$ या $3$ खोलता है (आपका कभी नहीं)। उनमें से दो में कार
दिख जाती है (कार 2 पर/खोला 2, कार 3 पर/खोला 3): वे हटा दिए जाते हैं। बचे
हुए चार में --- कार 1/खोला 2, कार 1/खोला 3, कार 2/खोला 3, कार 3/खोला 2 ---
बदलने पर अंतिम दो में जीत होती है: $\frac24 = \frac12$। अनजान संचालक के साथ 50--50
\emph{सही ही} है: विरोधाभास पूरा का पूरा संचालक के ज्ञान में बसता है, जो
परिदृश्यों को असमान रूप से छानता है।

\textbf{10.} आपके चुनाव में कार होने की \omterm{def:g10:proba:distribution}{प्रायिकता} $\frac{1}{100}$ है; जानकार संचालक की
$98$ खुलाइयाँ बची हुई $\frac{99}{100}$ को उस एक बंद दरवाज़े पर कीप की तरह उड़ेल देती
हैं। बदलिए, और $100$ में $99$ बार जीतिए।

\textbf{11.} तीन पत्ते (एक इक्का, दो जोकर), एक मित्र जो झाँककर देख लेता है
और आपके न चुने हुए दो में से हमेशा एक जोकर हटा देता है; $300$ दौर खेलिए,
हमेशा बदलिए, और जीत दर्ज कीजिए। अपेक्षित \omterm{def:g10:stats:series}{आपेक्षिक बारंबारता} $\frac23 \approx 0.667$, और
उतार-चढ़ाव $\pm\frac{1}{\sqrt{300}} \approx 0.058$: लगभग $61\,\%$ और $72\,\%$ के बीच बदलने पर जीत की अपेक्षा
कीजिए --- किसी भी ``50--50'' से कोसों दूर।

\textbf{12.} समसंभावी रचनाएँ (बड़ा--छोटा): BB, BG, GB, GG। ``कम से कम एक
लड़का'' BB, BG, GB रहने देता है: $P(\text{BB}) = \frac13$।

\textbf{13.} ``बड़ी संतान लड़का है'' BB, BG रहने देता है: $P(\text{BB}) = \frac12$। भिन्न सूचना
भिन्न परिदृश्य-समुच्चय छानती है --- ठीक वैसे ही जैसे दोनों मॉन्टी ने किया:
आप क्या जानते हैं यह मायने रखता है, और \emph{कैसे} जाना यह भी उतना ही।

\textbf{14.} तीनों एक जैसे: HHH, TTT: $\frac28 = \frac14$। दालम्बेर की त्रुटि: उनकी दोनों
``स्थितियाँ'' समसंभावी नहीं हैं --- ``तीनों एक जैसे नहीं'' आठ में से छह
परिणामों को एक गठरी में बाँध देती है। \Cref{prop:g10:proba:uniform} केवल समसंभावी परिणामों पर गिनती
करता है।

\textbf{15.} प्रति खेल औसत जीत: $6 \times \frac16 + 3
\times \frac16 = 1.50$ यूरो, जबकि दाँव $2$ यूरो का है:
प्रति खेल औसतन $0.50$ की हानि, और $600$ खेलों में लगभग $300$ यूरो। न्यायसंगत
प्रीमियम औसत भुगतान के बराबर होता है: $1\,\% \times 10\,000 = 100$ यूरो; वास्तविक प्रीमियम में
ख़र्च और सुरक्षा का अंतर जुड़ जाता है --- बीमाकर्ता को बुरे वर्षों में भी
जीवित रहना है, केवल औसत वर्षों में नहीं।

\textbf{16.} \omterm{def:g10:numbers:interval}{अंतराल}: $0.5 \pm \frac{1}{\sqrt{100}} =
\intcc{0.4}{0.6}$। $55$ चित ($0.55$): भीतर, कुछ ख़ास नहीं। $70$
चित ($0.70$): बहुत बाहर --- सिक्के पर संदेह कीजिए।

\textbf{17.} $\frac{1}{\sqrt{1000}} \approx 0.032$: सचमुच $50\,\%$ वाला उम्मीदवार प्रायः $\intcc{0.468}{0.532}$ में सर्वेक्षित
होता --- और $0.52$ उसके भीतर आराम से बैठा है। सर्वेक्षण बढ़त का \emph{संकेत}
देता है पर सिद्ध कुछ नहीं करता: अंतर उसे निगल जाता है।

\textbf{18.} पहला \omterm{def:g10:proba:sample}{प्रतिदर्श}: \omterm{def:g10:numbers:interval}{अंतराल} $0.04 \pm \frac{1}{20} =
\intcc{0}{0.09}$; प्रेक्षित $0.07$ भीतर है: कोई
निर्णय नहीं। दूसरा: $0.04 \pm \frac{1}{50} = \intcc{0.02}{0.06}$; प्रेक्षित $0.07$ बाहर पड़ता है: अब दावा गंभीर
संदेह में है। वही \omterm{def:g10:stats:series}{आपेक्षिक बारंबारता}, पर पैना यंत्र: यथार्थता $\sqrt n$ की तरह
बढ़ती है।

\textbf{19.} उतार-चढ़ाव $\frac{1}{\sqrt n}$ की तरह सिकुड़ता है: \omterm{def:g10:proba:sample}{प्रतिदर्श} चौगुना कीजिए,
शोर आधा हो जाएगा। अभिनति --- ग़लत भीड़ से पूछना --- किसी भी $n$ पर बची
रहती है: बीस लाख ग़लत चुने गए मतपत्र ग़लत ही रहे (माध्यमिक विद्यालय खंड की
``आँकड़े कैसे झूठ बोलते हैं'' वाली सप्ताहांत समस्या)।

\textbf{20.} (क) गिनने से पहले जाँचिए कि परिणाम समसंभावी हैं --- दालम्बेर
यही भूल गए। (ख) \omterm{met:g10:proba:tree}{वृक्ष} शाखाओं के अनुदिश गुणा करते हैं; ``कम से कम एक''
\omterm{def:g10:proba:operations}{पूरक} के आगे हथियार डाल देता है। (ग) सूचना परिदृश्य छानती है, और उसका
\emph{स्रोत} छन्नी बदल देता है --- जानकार मॉन्टी बनाम गिरा हुआ मॉन्टी,
बड़ा-लड़का बनाम कोई-एक-लड़का। (घ) आपेक्षिक बारंबारताएँ $\frac{1}{\sqrt n}$ की चाल से
\omterm{def:g10:proba:distribution}{प्रायिकताओं} के पास पहुँचती हैं, इसलिए धैर्य से किया गया अनुकरण हर बहस
सुलझा देता है --- उस बहस को भी जिसमें दस हज़ार नाराज़ पत्र-लेखक हों।
\end{solution}
